\section{Limits of Individual-Source Methods}
\label{sec:limits_individual}

The preceding chapters searched for dark matter in individual, resolved gamma-ray targets.
We now step back and assess the broader limitations of that paradigm, before motivating the transition to population-level methods that forms the core of Part III.

\subsection{The Resolved-Source Paradigm and Its Constraints}
\label{sec:6.1.1}

The Galactic Center Excess, first reported by Goodenough and Hooper \cite{Goodenough:2009gk} and subsequently confirmed by independent analyses \cite{Hooper:2010mq,Daylan:2014rsa,Calore:2014xka}, remains the most prominent anomaly in \textit{Fermi}-LAT data. As discussed in Chapter~\ref{ch:4} (Section~\ref{sec:4.1}), disentangling a potential dark matter contribution from astrophysical explanations has proven so far intractable.
The difficulty is not a lack of signal strength but an excess of competing models: the morphology of the excess is equally well described by an NFW-squared dark matter profile and by a population of unresolved millisecond pulsars tracing the stellar bulge \cite{Bartels:2015aea,Lee:2015fea}. The interstellar emission models introduce further degeneracies, as different templates for the Galactic diffuse emission shift the inferred GCE spectrum and morphology \cite{Calore:2014xka}. Even the non-Poissonian template fitting technique, initially presented as a way to distinguish point-source from diffuse origins, was shown to be sensitive to mismodeling artifacts \cite{Leane:2019xiy,Leane:2020nmi}.
More recently, simulation-based inference approaches incorporating energy information have reopened the question.
As discussed in Section~\ref{sec:4.3.3}, the CNN-based analysis of List et al.~\cite{List:2025qbx} pushes the inferred source-count distribution to fluxes so faint that the GCE becomes statistically indistinguishable from the smooth, Poissonian emission expected from dark matter annihilation.
The origin of the excess thus remains an open problem, with neither the dark matter nor the millisecond pulsar hypothesis decisively favored by current data.
The Galactic Center therefore illustrates a general principle: when the target region contains multiple bright, partially degenerate emission components, attributing any residual to a specific physical origin becomes a problem of systematic uncertainties rather than statistical sensitivity.

Chapter~\ref{ch:5} pursued an alternative strategy, searching for dark matter subhalos among the 2428 unassociated point sources in the 4FGL-DR4 catalog \cite{Fermi-LAT:2022byn,Ballet:2023qzs} at Galactic latitude $|b| > 10^\circ$ (Section~\ref{sec:5.4}).
The environment is far cleaner: high-latitude sources are not contaminated by the Galactic diffuse emission that plagues the GCE.
However, our quantification analysis found no statistically significant dark matter subhalo component (Section~\ref{sec:halos:results}).
This null result reflects two structural constraints.
First, only the most massive and nearby subhalos produce fluxes above the \textit{Fermi}-LAT catalog threshold of $\mathrm{TS} > 25$; the vast majority of the predicted subhalo population \cite{Springel:2008cc} lies well below it.
Second, the look-elsewhere penalty is severe: any individual candidate extracted from a catalog of thousands of unassociated sources requires evidence well beyond the $5\sigma$ threshold conventionally demanded for discovery.

These two case studies expose a shared limitation.
Both the GCE analysis and the subhalo search operate on sources that have already been detected or on emission that is already visible.
They are structurally blind to the population of sources below the detection threshold.
Yet this sub-threshold population is far from negligible: the very existence of the isotropic diffuse gamma-ray background \cite{Fermi-LAT:2014ryh} implies a substantial flux contribution from sources too faint to be individually resolved.
The challenge is not to look harder at individual sources, but to develop methods that extract population-level information from the collective photon statistics of the unresolved sky.

\subsection{The Population-Level Alternative}
\label{sec:6.1.2}

The limitations of individual source identification motivate a conceptual shift.
Rather than asking whether any specific source is a dark matter subhalo, we can ask a different question: what is the statistical distribution of all sources, including those too faint to detect individually?

Below the formal \textit{Fermi}-LAT detection threshold lies a vast population of unresolved sources.
Their cumulative emission forms a significant fraction of the unresolved gamma-ray background (UGRB) \cite{DGRB-review,Fermi-LAT:2014ryh}.
While these sources are invisible individually, their collective photon statistics encode recoverable information about the underlying population.
This population-level approach complements the strategies explored in Part II: resolved-source methods probe the bright tail of the population, whereas statistical methods probe the faint end.
Together, they provide a complete census.

Chapters~\ref{ch:6} and \ref{ch:7} focus on characterizing this general unresolved source population --- primarily astrophysical objects such as blazars \cite{Ajello:2015mfa,DiMauro:2017ing}, millisecond pulsars \cite{Manconi:2019ynl}, and star-forming galaxies \cite{Linden:2016fdd}.
A thorough census of the astrophysical population is a prerequisite for any dedicated dark matter search: only after accounting for the known source classes can a residual exotic component be meaningfully constrained.
The search for a collective dark matter signal in the unresolved gamma-ray sky is deferred to Chapter~\ref{ch:8}, which employs angular cross-correlations with tracers of large-scale structure, a technique complementary to the population statistics developed here.

The primary observable for characterizing the sub-threshold source population is the source-count distribution, denoted $dN/dS$.
We define this distribution, discuss its physical interpretation, and review analytical approaches to measuring it in the following section.
